% Generated by build/sync_qft_soln.py from committed sources only.
% Source repository: https://github.com/Luca-Yucheng-Jin/QFT-soln
% Source commit: 07d7fe95fd5cf5b26d38a16ee535a04a925277b2
\section{$SU(N)$ Yang--Mills Theory and Its Coupling to Scalar Fields (PSI QFT II, Tutorial 8)}
\markboth{\thesection\quad $SU(N)$ YANG--MILLS WITH SCALARS (PSI QFT II, TUTORIAL 8)}{}

\noindent\textit{Question prompts paraphrased from the official
\href{https://psi-online.perimeterinstitute.ca/courses/take/quantum-field-theory-ii-student/pdfs/20776405-module-7-tutorial}{Perimeter PSI Online sheet};
the equations and notation follow the source.}

\subsection{Coupling to Scalar Fields}

Let $\phi_i(x)$ be a scalar field in a representation $r$ of $SU(N)$, with
$i=1,\ldots,\dim r$.

\begin{enumerate}[label=\alph*)]
    \item Determine the dimension of the group $SU(N)$.
    \begin{solution}
        An $N\times N$ complex matrix has $2N^2$ real dimensions. The hermitian condition, $U^*_{ik}U_{jk} = \delta_{ij}$, is $N^2$ constraints. In addition, $\det U =1$ gives an extra constraint. Therefore, $\dim SU(N) = N^2 -1$.
    \end{solution}

    \item Determine the dimension of the fundamental representation of $SU(N)$.
    \begin{solution}
        The fundamental representation of $SU(N)$ is $N \times N$ complex matrices, which has dimension $N$.
    \end{solution}

    \item The adjoint action is defined by
    \begin{equation}
        \operatorname{Ad}_g(T^a)\equiv gT^ag^{-1}.
    \end{equation}
    Thus there is a $\dim G\times\dim G$ matrix $C(g)$ such that
    \begin{equation}
        gT^ag^{-1}=C^{ab}(g)T^b.
    \end{equation}
    Expand $g$ and $C(g)$ infinitesimally and use
    \begin{equation}
        [T_i,T_j]=if_{ijk}T_k
    \end{equation}
    to derive the adjoint generators
    \begin{equation}
        (T_i)_{jk}=if_{ijk}.
    \end{equation}
    \begin{solution}
        \begin{equation}
            \begin{aligned}
                \operatorname{Ad}(g)Y = (I + \epsilon X)Y(I-\epsilon X) = I + \epsilon[X,Y] = (I + \epsilon \operatorname{ad}(X))Y \implies \operatorname{ad}(X)Y = [X,Y].
            \end{aligned}
        \end{equation}
        Taking $X = T^a, Y=T^b$, we have
        \begin{equation}
            \operatorname{ad}(T^a)^{bc}T^c = if^{abc}T^c \implies \operatorname{ad}(T^a)^{bc} =if^{abc}.
        \end{equation}
    \end{solution}

    \item Determine the dimension of the adjoint representation of $SU(N)$.
    \begin{solution}
        As we have seen above, the adjoint representation of the Lie algebra is a $\dim G \times \dim G$ matrix. Therefore,
        \begin{equation}
            \dim \operatorname{adj}(su(N)) = \dim SU(N).
        \end{equation}
    \end{solution}
\end{enumerate}

Now consider the local transformation
\begin{equation}
    \phi_i(x)\longmapsto g_{ij}(x)\phi_j(x),
    \qquad
    g(x)=\exp\!\bigl[i\alpha_m(x)T_r^m\bigr],
\end{equation}
where $T_r^m$ are the $su(N)$ generators in the representation $r$ and the
$\alpha_m(x)$ are spacetime-dependent group parameters. Require the covariant
derivative to transform as
\begin{equation}
    D_\mu\phi_i(x)\longmapsto g_{ij}(x)D_\mu\phi_j(x),
\end{equation}
and define, with matrix indices suppressed,
\begin{equation}
    D_\mu\phi=\partial_\mu\phi-ieA_\mu\phi,
    \qquad
    A_\mu(x)=A_\mu^m(x)T_r^m.
\end{equation}

\begin{enumerate}[label=\alph*),resume]
    \item Derive the finite gauge-transformation law of $A_\mu(x)$ required by
    the covariance of $D_\mu\phi$.
    \begin{solution}
        \begin{equation}
            \begin{aligned}
                D_\mu \phi \to D_\mu' g \phi = (\partial_\mu g)\phi + g \partial_\mu \phi  - ie A_\mu' g\phi = g(\partial_\mu \phi - ie A_\mu \phi),
        \end{aligned}
        \end{equation}
        which implies
        \begin{equation}
            A_\mu' = g A_\mu g^{-1} - \frac{i}{e}(\partial_\mu g) g^{-1}.
        \end{equation}
    \end{solution}

    \item For infinitesimal $\alpha_m(x)$, show that
    \begin{equation}
        \delta A_\mu^a=\frac{1}{e}D_\mu\alpha_a,
        \qquad
        D_\mu\alpha=\partial_\mu\alpha+ie[\alpha,A_\mu].
    \end{equation}
    \begin{solution}
        \begin{equation}
            \begin{aligned}
                A_\mu^aR(T^a) & \to (I + i \alpha_bR(T^b))A_\mu^a R(T^a)(I - i \alpha_bR(T^b)) + \frac{1}{e}(\partial_\mu \alpha_b)R(T^b)(I + i \alpha_a R(T^a))\\
                &=A_\mu^aR(T^a) + i\alpha_b[R(T^b),R(T^a)]A_\mu^a+\frac{1}{e}(\partial_\mu \alpha_b)R(T^b)\\
                &=A_\mu^aR(T^a) + \frac{1}{e}D_\mu\alpha,
            \end{aligned}
        \end{equation}
        from which we can recognise $\alpha$ as the adjoint representation of $su(N)$.
    \end{solution}

    \item Prove that the Lagrangian
    \begin{equation}
        \mathcal{L}
        =-(D_\mu\phi)^\dagger D^\mu\phi
        -m^2\phi^\dagger\phi
        -\frac{\lambda}{2}(\phi^\dagger\phi)^2
    \end{equation}
    is invariant under the local scalar transformation and the gauge-field
    transformation found above.
    \begin{solution}
        It can be easily seen through the transformation of $\phi$ and $D_\mu \phi$.
    \end{solution}

    \item Extract every Feynman rule from this Lagrangian that contains at least
    one scalar line.
    \begin{solution}
        Expanding the covariant derivative, we have
        \begin{equation}
            |D_\mu\phi|^2  = |\partial_\mu\phi|^2 + ieA_\mu(\phi^\dagger\partial_\mu \phi - \phi \partial_\mu \phi^\dagger) + e^2A_\mu^2|\phi|^2,
        \end{equation}
        and the Feynman rules can be read off.
    \end{solution}

    \item Apply Noether's theorem to obtain the conserved currents.
    \begin{solution}
        \begin{equation}
            j^a_\mu = \frac{\partial \mathcal{L}}{\partial(\partial^\mu \phi)}\delta\phi^a  +  \frac{\partial \mathcal{L}}{\partial(\partial^\mu \phi^\dagger)}\delta\phi^{\dagger a}= i (D_\mu\phi)^\dagger R(T^a)\phi - i\phi^\dagger R(T^a)(D_\mu \phi).
        \end{equation}
    \end{solution}
\end{enumerate}

\subsection{Zero Modes of the Yang--Mills Lagrangian}

Take the unfixed Yang--Mills Lagrangian
\begin{equation}
    \mathcal{L}
    =-\frac{1}{2g^2}\operatorname{Tr}(F^{\mu\nu}F_{\mu\nu}),
\end{equation}
with the trace over adjoint indices and arbitrary $N$ in $SU(N)$.

\begin{enumerate}[label=\alph*)]
    \item Expand the $F^2$ term, isolate the quadratic terms containing two
    derivatives, and write them as
    \begin{equation}
        \int d^4x\,A_\mu^a(x)D_x^{\mu\nu}A_\nu^a(x).
    \end{equation}
    \begin{solution}
        \begin{equation}
            \begin{aligned}
                S &= -\frac{1}{4g}\int d^dx F_{\mu\nu}^a F^{a\mu\nu}\\
                &=-\frac{1}{2g}\int d^dx A_\nu(-\partial^2\eta^{\mu\nu} + \partial^\mu\partial^\nu)A_\mu + \cdots,
            \end{aligned}
        \end{equation}
        where the interacting terms are suppressed.
    \end{solution}

    \item Show that there is a nonzero field $B_\mu$ satisfying
    \begin{equation}
        D_x^{\mu\nu}B_\nu(x)=0.
    \end{equation}
    \textit{Hint: regard $D^{\mu\nu}$ as an integral kernel and inspect its
    momentum-space representation.}
    Explain why this zero mode obstructs inversion of the quadratic operator
    before gauge fixing.
    \begin{solution}
        It is easy to see that any field of the form $\partial_\mu \omega$ is the kernel of the operator $D^{\mu\nu}_a$. Therefore, $D^{\mu\nu}_a$ is clearly not invertible and hence we need the Faddeev-Popov procedure.
    \end{solution}
\end{enumerate}
